\chapter[\emj{🧮}~Matemáticas: GCD, Criba y Exponenciación Rápida]{\emj{🧮}~Algoritmos Matemáticos Esenciales}

\begin{dsaquees}

Un puñado de trucos matemáticos que aparecen escondidos en muchísimos
problemas 🧮:

\begin{itemize}
\item GCD (Euclides): el máximo común divisor de dos números, en un
      parpadeo.
\item Criba de Eratóstenes: encontrar TODOS los números primos hasta N
      rapidísimo.
\item Exponenciación rápida: calcular $a^b$ en $O(\log b)$ en vez de
      multiplicar $b$ veces.
\item Aritmética modular: trabajar con residuos para no desbordar con
      números gigantes.
\end{itemize}

\end{dsaquees}

\begin{dsanino}

Cuatro herramientas de matemáticas mágicas 🎩:

GCD: para saber el pedazo más grande con el que puedes medir dos cuerdas
🪢 de distinto largo sin que sobre. Vas restando la chica de la grande (o
usando el residuo) hasta que quedan iguales: ese es el GCD.

Criba: quieres una lista de números primos (los que solo se dividen entre
1 y ellos mismos). Tachas todos los múltiplos de 2, luego de 3, luego de
5\ldots y los que quedan sin tachar son primos.

Exponenciación rápida: para calcular 2 elevado a 16, en vez de multiplicar
2 dieciséis veces, elevas al cuadrado: 2, 4, 16, 256\ldots ¡en 4 pasos!
Cada vez duplicas el exponente. Menos trabajo, mismo resultado.

\end{dsanino}

\begin{dsaotraforma}

\textbf{GCD (algoritmo de Euclides)}: $\gcd(a, b) = \gcd(b, a \bmod b)$
hasta que $b = 0$. $O(\log \min(a,b))$. De ahí sale el LCM:
$\mathrm{lcm}(a,b) = a / \gcd(a,b) \cdot b$.

\textbf{Criba de Eratóstenes}: marca como no-primos los múltiplos de cada
primo desde 2. Genera todos los primos hasta $N$ en
$O(N \log \log N)$.

\textbf{Exponenciación rápida (binary exponentiation)}: usa la
representación binaria del exponente. Si $b$ es par,
$a^b = (a^{b/2})^2$; si es impar, $a^b = a \cdot a^{b-1}$. $O(\log b)$.
Combinada con módulo evita overflow: base de la exponenciación modular
usada en criptografía y hashing.

\textbf{Aritmética modular}: $(x + y) \bmod m$, $(x \cdot y) \bmod m$
mantienen los números acotados. Regla clave: aplica el módulo tras cada
operación para no desbordar.

\end{dsaotraforma}

\begin{dsadiagrama}
\begin{verbatim}
EUCLIDES  gcd(48, 18):
  gcd(48,18) -> 48 mod 18 = 12 -> gcd(18,12)
  gcd(18,12) -> 18 mod 12 = 6  -> gcd(12,6)
  gcd(12,6)  -> 12 mod 6  = 0  -> gcd(6,0) = 6
  GCD = 6

CRIBA hasta 20 (X = tachado, no primo):
  2  3  4X 5  6X 7  8X 9X 10X 11 12X 13 14X 15X 16X 17 18X 19 20X
  Primos: 2 3 5 7 11 13 17 19

EXP RÁPIDA  2^10:
  10 = 1010 binario
  2^10 = 2^8 * 2^2 = 256 * 4 = 1024  (4 cuadrados, no 10 multiplic.)
\end{verbatim}
\end{dsadiagrama}

\begin{lstlisting}[language=C++]
GCD (Euclides)
- gcd(a, b): si b == 0 return a; si no gcd(b, a % b).

EXPONENCIACIÓN RÁPIDA (con módulo)
- Mientras el exponente > 0: si es impar, multiplica el resultado por la
  base; eleva la base al cuadrado; divide el exponente entre 2.

📝 PSEUDOCÓDIGO (exponenciación modular)
──────────────────────────────────────────────────────────────

función powMod(a, b, m):
    resultado = 1; a = a mod m
    MIENTRAS b > 0:
        SI b es impar: resultado = resultado * a mod m
        a = a * a mod m
        b = b / 2            // desplaza un bit
    return resultado

💻 CÓDIGO C++
──────────────────────────────────────────────────────────────

#include <vector>
using namespace std;

// GCD y LCM
long long gcd(long long a, long long b) {
    return b == 0 ? a : gcd(b, a % b);
}
long long lcm(long long a, long long b) {
    return a / gcd(a, b) * b;         // divide primero: evita overflow
}

// Exponenciación modular rápida -> O(log b)
long long powMod(long long a, long long b, long long m) {
    long long res = 1; a %= m;
    while (b > 0) {
        if (b & 1) res = res * a % m;  // bit impar
        a = a * a % m;                 // eleva al cuadrado
        b >>= 1;                       // siguiente bit
    }
    return res;
}

// Criba de Eratóstenes -> primos hasta n
vector<bool> criba(int n) {
    vector<bool> esPrimo(n + 1, true);
    esPrimo[0] = esPrimo[1] = false;
    for (long long i = 2; i * i <= n; i++)
        if (esPrimo[i])
            for (long long j = i * i; j <= n; j += i)  // desde i*i
                esPrimo[j] = false;
    return esPrimo;
}

⏱️ COMPLEJIDAD (Big-O)
──────────────────────────────────────────────────────────────

  Algoritmo         │ Tiempo
  ──────────────────┼──────────────────
  GCD (Euclides)    │ O(log min(a,b))
  Exp. rápida       │ O(log b)
  Criba             │ O(n log log n)

* La criba usa O(n) memoria (un bool por número hasta n).
\end{lstlisting}

\begin{dsaentrevistas}

✅ powMod es la base de todo: hashing polinómico, inverso modular (por
   Fermat: $a^{m-2} \bmod m$ cuando $m$ es primo), y criptografía. Tenlo
   memorizado.

💡 Criba desde i·i: al marcar múltiplos de $i$, empieza en $i^2$ (los
   menores ya fueron marcados por primos más chicos). Pequeña
   optimización que preguntan.

⚠️ Overflow: en \verb|a * a % m| con long long y $m$ grande puedes
   desbordar. Si $m \sim 10^{18}$, usa \verb|__int128| o multiplicación
   modular segura.

🔑 GCD para simplificar: muchos problemas de fracciones, rotaciones de
   arreglos ("rotate array by k" usa \verb|gcd(n,k)| ciclos) y geometría
   de rejilla se apoyan en el GCD.

\end{dsaentrevistas}

\begin{dsaresumen}

• GCD (Euclides): gcd(a,b) = gcd(b, a mod b); O(log min).
• LCM = a / gcd(a,b) · b (divide primero para no desbordar).
• Criba: tacha múltiplos desde i·i; primos hasta n en O(n log log n).
• Exponenciación rápida: eleva al cuadrado; O(log b).
• powMod evita overflow y es base de hashing e inverso modular.
• Aritmética modular: aplica el módulo tras cada operación.
════════════════════════════════════════════════════════════════

\end{dsaresumen}
